\section{Protokol Pertukaran Kunci Diffie-Hellman}
Dipublikasikan oleh Whitfield Diffie dan Martin Hellman pada tahun 1976 \citep{diffie1976}, protokol ini memungkinkan dua pihak (misalnya Alice dan Bob) untuk menyepakati sebuah kunci rahasia bersama melalui saluran publik yang tidak aman.

\subsection{Parameter Publik}
Sebelum memulai pertukaran, Alice dan Bob menyepakati dua parameter publik yang tidak perlu dirahasiakan:
\begin{enumerate}
    \item Sebuah bilangan prima besar, $p$.
    \item Sebuah bilangan bulat $g$, sedemikian sehingga $g$ adalah akar primitif dari $p$ (di mana $g < p$).
\end{enumerate}

\subsection{Prosedur Pertukaran Kunci}
Langkah-langkah perhitungan Diffie-Hellman adalah sebagai berikut:
\begin{enumerate}
    \item \textbf{Pembangkitan Kunci Privat}:
    \begin{itemize}
        \item Alice memilih bilangan bulat acak $a$ sebagai kunci privatnya.
        \item Bob memilih bilangan bulat acak $b$ sebagai kunci privatnya.
    \end{itemize}
    \item \textbf{Perhitungan dan Pertukaran Kunci Publik}:
    \begin{itemize}
        \item Alice menghitung $A = g^a \bmod p$ dan mengirimkan $A$ kepada Bob.
        \item Bob menghitung $B = g^b \bmod p$ dan mengirimkan $B$ kepada Alice.
    \end{itemize}
    \item \textbf{Penghitungan Kunci Rahasia Bersama}:
    \begin{itemize}
        \item Alice menghitung $K = B^a \bmod p$.
        \item Bob menghitung $K = A^b \bmod p$.
    \end{itemize}
\end{enumerate}

\textbf{Pembuktian Matematis}:
Kunci yang dihasilkan oleh Alice dan Bob adalah sama karena:
\[ K = B^a \bmod p = (g^b)^a \bmod p = g^{ba} \bmod p = (g^a)^b \bmod p = A^b \bmod p \]

% Pertukaran kunci Diffie-Hellman sebagai diagram urutan.
% Dirujuk dari section-02.tex dengan % Pertukaran kunci Diffie-Hellman sebagai diagram urutan.
% Dirujuk dari section-02.tex dengan % Pertukaran kunci Diffie-Hellman sebagai diagram urutan.
% Dirujuk dari section-02.tex dengan \input{figures/alur-dh}.
\begin{figure}[htbp]
    \centering
    \begin{tikzpicture}[
        panah/.style={-{Latex[length=2.4mm]}, thick},
        label/.style={font=\scriptsize, align=center, fill=white, inner sep=2pt},
        % Anotasi di sisi lifeline: `fill=white` menutup garis putus-putus
        % agar teks tidak tertimpa.
        sisi/.style={font=\scriptsize, align=center, fill=white, inner sep=2pt},
        pihak/.style={draw, rounded corners, minimum width=2.8cm, minimum height=0.7cm},
    ]
        \node[pihak, fill=blue!6]  (a) at (0,0) {\small Alice};
        \node[pihak, fill=green!10] (b) at (7,0) {\small Bob};
        \draw[dashed] (a.south) -- ++(0,-4.75);
        \draw[dashed] (b.south) -- ++(0,-4.75);

        \node[sisi] at (0,-1.1) {parameter publik\\ $p$ dan $g$};
        \node[sisi] at (7,-1.1) {parameter publik\\ $p$ dan $g$};
        \node[sisi] at (0,-2.3) {pilih $a$ rahasia};
        \node[sisi] at (7,-2.3) {pilih $b$ rahasia};

        \draw[panah] (0,-3.3) -- node[label] {$A = g^{a} \bmod p$} (7,-3.3);
        \draw[panah] (7,-4.2) -- node[label] {$B = g^{b} \bmod p$} (0,-4.2);

        \node[sisi] at (0,-5.0) {$K = B^{a} \bmod p$};
        \node[sisi] at (7,-5.0) {$K = A^{b} \bmod p$};
    \end{tikzpicture}
    \caption{Pertukaran kunci Diffie-Hellman: kedua pihak menyepakati kunci rahasia $K$ yang sama walau hanya bertukar $A$ dan $B$ di saluran publik}
    \label{fig:alur-dh}
\end{figure}
.
\begin{figure}[htbp]
    \centering
    \begin{tikzpicture}[
        panah/.style={-{Latex[length=2.4mm]}, thick},
        label/.style={font=\scriptsize, align=center, fill=white, inner sep=2pt},
        % Anotasi di sisi lifeline: `fill=white` menutup garis putus-putus
        % agar teks tidak tertimpa.
        sisi/.style={font=\scriptsize, align=center, fill=white, inner sep=2pt},
        pihak/.style={draw, rounded corners, minimum width=2.8cm, minimum height=0.7cm},
    ]
        \node[pihak, fill=blue!6]  (a) at (0,0) {\small Alice};
        \node[pihak, fill=green!10] (b) at (7,0) {\small Bob};
        \draw[dashed] (a.south) -- ++(0,-4.75);
        \draw[dashed] (b.south) -- ++(0,-4.75);

        \node[sisi] at (0,-1.1) {parameter publik\\ $p$ dan $g$};
        \node[sisi] at (7,-1.1) {parameter publik\\ $p$ dan $g$};
        \node[sisi] at (0,-2.3) {pilih $a$ rahasia};
        \node[sisi] at (7,-2.3) {pilih $b$ rahasia};

        \draw[panah] (0,-3.3) -- node[label] {$A = g^{a} \bmod p$} (7,-3.3);
        \draw[panah] (7,-4.2) -- node[label] {$B = g^{b} \bmod p$} (0,-4.2);

        \node[sisi] at (0,-5.0) {$K = B^{a} \bmod p$};
        \node[sisi] at (7,-5.0) {$K = A^{b} \bmod p$};
    \end{tikzpicture}
    \caption{Pertukaran kunci Diffie-Hellman: kedua pihak menyepakati kunci rahasia $K$ yang sama walau hanya bertukar $A$ dan $B$ di saluran publik}
    \label{fig:alur-dh}
\end{figure}
.
\begin{figure}[htbp]
    \centering
    \begin{tikzpicture}[
        panah/.style={-{Latex[length=2.4mm]}, thick},
        label/.style={font=\scriptsize, align=center, fill=white, inner sep=2pt},
        % Anotasi di sisi lifeline: `fill=white` menutup garis putus-putus
        % agar teks tidak tertimpa.
        sisi/.style={font=\scriptsize, align=center, fill=white, inner sep=2pt},
        pihak/.style={draw, rounded corners, minimum width=2.8cm, minimum height=0.7cm},
    ]
        \node[pihak, fill=blue!6]  (a) at (0,0) {\small Alice};
        \node[pihak, fill=green!10] (b) at (7,0) {\small Bob};
        \draw[dashed] (a.south) -- ++(0,-4.75);
        \draw[dashed] (b.south) -- ++(0,-4.75);

        \node[sisi] at (0,-1.1) {parameter publik\\ $p$ dan $g$};
        \node[sisi] at (7,-1.1) {parameter publik\\ $p$ dan $g$};
        \node[sisi] at (0,-2.3) {pilih $a$ rahasia};
        \node[sisi] at (7,-2.3) {pilih $b$ rahasia};

        \draw[panah] (0,-3.3) -- node[label] {$A = g^{a} \bmod p$} (7,-3.3);
        \draw[panah] (7,-4.2) -- node[label] {$B = g^{b} \bmod p$} (0,-4.2);

        \node[sisi] at (0,-5.0) {$K = B^{a} \bmod p$};
        \node[sisi] at (7,-5.0) {$K = A^{b} \bmod p$};
    \end{tikzpicture}
    \caption{Pertukaran kunci Diffie-Hellman: kedua pihak menyepakati kunci rahasia $K$ yang sama walau hanya bertukar $A$ dan $B$ di saluran publik}
    \label{fig:alur-dh}
\end{figure}


